\chapter{The only lost cause (Can You Face Your True Feelings)}
Last episode, we saw how Kyubey denies the agency of the magical girls by restricting their access to information. This episode opens by continuing that thought, as Kyubey demonstrates to the audience and Sayaka how little he cares about her. As part of a demonstration of his claim that the removal of the magical girls' souls is beneficial to them, he uses Sayaka's Soul Gem to torture her, demonstrating that the pain of the first blow in her fight with Kyoko would have crippled her utterly if not for the buffer provided by the gem. His total lack of interest in her agony as anything but a teaching tool, however, belies any claim by Kyubey to have the benefit of the magical girls in mind.

The repeated use of the word "zombie" to describe the magical girls (which is original to the Japanese text—Sayaka can be distinctly heard using it several times in the Japanese dialogue in the episode) is telling here. The entire point of a zombie is that it is shaped like a person, but actually a thing. In philosophy, a "zombie" is a creature that acts like a human but has no internal experience or life—for example, poking a zombie will cause it to say "ow," but it has no internal sense of pain \cite{ref50}. More familiarly, the zombies of movie fame are walking corpses, who can be fought and killed while technically being already dead. This allows the audience the visceral thrill of imagining fighting and killing other people, without having to worry about the morality of actually causing a human being to die \cite{ref51}. In other words, both the philosophical zombie and the movie zombie are extreme cases of objectification, in which a person's agency is stripped away leaving only a thing which can be used and abused with impunity. Arguably, the zombie myth even originates from feelings of being unfairly used, dehumanized, or stifled by an overwhelming and interloping alien power \cite{ref52}; certainly a reasonable description of Sayaka's current state and Kyubey's influence on her.

Objectification is rampant throughout the episode. Paralleled to Kyubey, who assumes he knows what is best for Sayaka and causes her intense suffering as a consequence, is Kyoko. Throughout Kyoko's story of how she became a magical girl, her family and the people around them are represented by dolls, puppets, and toys—like zombies, human-shaped objects that possess no agency. Kyoko made a wish that she believed was what her father wanted, but since she did not consult him, she got it wrong, and as a consequence brought pain not only to her father but to her entire family and ultimately herself. By failing to communicate openly with her father, and simply assuming she knew what was best for him, she treated him as the object of her observations, rather than as a subject capable of expressing his own needs and wishes. As much as unspoken communication and indirect communication styles are common in Japan \cite{ref53}, this preference for indirectness ended up causing a misunderstanding of such proportions that it still haunts the seemingly-uncaring Kyoko.

However, Kyoko has learned entirely the wrong lessons from this experience. Rather than talking to Sayaka and trying to understand her, she assumes that she and Sayaka are the same, and that Sayaka's problem is that she wished for Kyousuke without understanding what Kyousuke wanted. Kyoko has come to reject empathy and human connection entirely, fixating on purely physical needs (namely, food) as a source of comfort while rejecting all human companionship and norms. This is unacceptable to Sayaka, because Sayaka's problem is ultimately not that she objectifies others.

Sayaka spoke with Kyousuke frequently while he was hospitalized; we saw him express quite clearly that he wanted his arm to heal, and that its inability to heal was a source of despair to him. Sayaka did not assume that she knew what Kyousuke wanted; she found that out directly from him. Rather, Sayaka's error was (as Mami hinted back in Episode 2) not understanding what she wanted—she wanted a healed and happy Kyousuke, yes, but specifically so that she could be with him. Sayaka's error, in other words, is that she is excessively self-sacrificing, in effect objectifying herself. 

Kyousuke's self-centered disinterest in Sayaka doesn't help her mental state at all, of course. He never thinks to question why she visited him so frequently in the hospital; he simply accepts this as normal and thus does not stop to consider whether Sayaka might appreciate being told when he leaves the hospital or that he is returning to school. He displays no interest in her inner life or motivations, and as such does not think to consider her feelings; he treats her as a secondary character in his life, rather than the main character of her own life—much as he himself is a secondary character in the show.

The cure for all the objectification going on in this episode, of course, is for characters to treat one another as agents possessed of unique, subjective experiences. Key to this is open communication, which unfortunately is in short supply throughout the episode. Characters mostly talk at each other, missing entirely the effects their words are having; the only real exception is the conversation between Madoka and Sayaka in which the latter breaks down, sobbing that she cannot approach Kyousuke romantically because of her altered physical state.

The trigger for that conversation is a demonstration of how difficult it can be to understand and appreciate others' subjectivity. Hitomi does everything right when she approaches Sayaka; nothing requires Hitomi to delay approaching Kyousuke or give Sayaka the first chance, but Hitomi does so anyway because she knows how Sayaka feels and doesn't want to hurt her. Unfortunately, because she doesn't know about everything else Sayaka is going through—and Sayaka understandably chooses not to tell her—she has no idea that by giving Sayaka this ultimatum she is triggering all of Sayaka's newly acquired body image issues. Hitomi has no way of knowing how she is hurting Sayaka, and likewise has no idea why Sayaka doesn't make her feelings known to Kyousuke; Hitomi thus has no reason to believe that Sayaka has any objections to or issues with Hitomi and Kyousuke dating.

Interestingly, however, it is not the loss of her love interest that most hurts Sayaka, but rather the brief moment during this conversation in which she regrets saving Hitomi from the witch in Episode 4. Sayaka is holding herself to a ridiculously high standard here, and thus failing to recognize that brief ugly impulses are a part of the human condition, an element of the internal life that does not necessarily translate into outward behavior. Instead, Sayaka takes this momentary viciousness as proof that she has truly become subhuman, that she is a "zombie" rather than a person with an unusual physical configuration.

Ultimately, this tendency of Sayaka to objectify herself culminates in deliberately numbing herself so that she can fight witches with no sense of pain, relying on her healing magic to repair the damage they do to her. Her sense of self-worth has plummeted to the point that she no longer cares about self-preservation and is no longer willing to accept help. It is only a matter of time, in her eyes, before she inevitably loses Kyousuke, and she feels that this is only right because she sees herself as having become a thing. At this point, despair and deep depression are all Sayaka sees in her future, and fighting witches the only purpose she has left.

The stage is now set for the culmination of the middle arc of Madoka Magica. The first arc ended with the show escaping from the constraints of the magical girl genre. This second arc will end with the genre's death.
